%% kapa: exam=1 round=21 year=2010 period=1 form=A subject=civil_law subject_ko=민법 seq=34 number=34
%% source: Q-Net 감정평가사 기출문제 (cst00309) artlSeq=5024181 file=extracted/1008/5024181/2010년 제21회 감정평가사 1차시험 기출문제/1교시.hwp
%% conversion: kordoc 4.24.51 -> tools/gen1.py (auto); review=proofread-sonnet; flags=none
\begin{question}{34}
소멸시효의 기산점에 관한 설명으로 옳은 것은? (다툼이 있으면 판례에 의함)
\begin{choices}
  \item 부작위를 목적으로 하는 채권의 소멸시효는 그 채권의 성립 시부터 진행한다.
  \item 동시이행의 항변권이 붙어 있는 채권의 경우, 항변권의 행사 시부터 소멸시효가 진행한다.
  \item 이행불능으로 인한 손해배상청구권의 소멸시효는 이행불능 시부터 진행한다.
  \item 장기간 입원치료를 받는 경우, 의사의 치료에 관한 개개의 채권은 특약이 없는 한 퇴원 시부터 진행한다.
  \item 정지조건부 권리는 권리가 성립한 때부터 시효가 진행한다.
\end{choices}
\end{question}
\answer{3}
